\chapter{双纽线动力学：感知—行动双循环与层级协同}
\label{chp:lemniscate_ch}

前两章分别建立了内部计算循环与外部交互事务，本章进一步回答两类循环怎样形成一个持续运行而又可检验的整体。我们把“双纽线”当作感知回路和行动回路共享边界接口的结构图，而不把图形相似性当作物理定律。研究对象包括智能体内部状态、环境状态、观测、控制、社会消息与层级指令；核心问题是这些异质对象如何通过带类型接口交换、怎样形成快慢反馈、何时保持稳定，以及能力不对称时如何避免控制关系被误写成必然支配。由此，我们既保留内部模型与现实约束相互校正的原问题，也把几何同构、相位纠缠、维数压制和“现在奇点”等强断言改写为可观测、可反驳的计算命题。

\section{双域系统：内部状态、环境过程与边界接口}
\label{sec:section_2656}

我们先定义双循环两侧的对象。智能体内部状态记为 $x(t)\in\mathcal{X}\subseteq\mathbb{R}^{n_x}$，其中可以包含工作状态 $h(t)$、情景记录 $E$ 的当前索引、结构参数 $\Theta$、任务目标 $g(t)$ 与资源账本 $b(t)$；环境状态记为 $z(t)\in\mathcal{Z}$，它可以是连续物理量、离散对象关系或二者组成的混合状态。内部状态不是环境的副本，环境状态也不因进入模型而变成语义向量。观测接口 $O_t:\mathcal{Z}\times\mathcal{N}\rightarrow\mathcal{Y}$ 把环境与传感噪声映射为观测 $y_t$，行动接口 $A_t:\mathcal{X}\times\mathcal{C}\rightarrow\mathcal{U}$ 把内部控制候选和接口约束映射为可执行控制 $u_t$。若 $y_t$ 具有伏特、像素、牛顿或离散事件等单位，编码器必须声明怎样把这些量变成无量纲内部坐标；若 $u_t$ 最终驱动真实装置，执行器必须声明量纲恢复、饱和、权限和安全限制。

内部表示仍由两类互补结构组成。局部结构网络 $G_t=(V_t,R_t)$ 保存对象身份、角色、来源、依赖与版本；全局分布表示 $q_t\in\mathcal{P}(\Omega_t)$ 或 $v_t\in\mathbb{R}^{d}$ 表达候选、语境和不确定性。两者通过绑定映射 $B_t:V_t\times\Omega_t\rightarrow[0,1]$、耦合算子和任务读出联合，而不是分别对应某个脑区或外界物体。环境侧也可能有图结构和连续变量，但内部图与环境图只在声明的查询族上比较。我们用任务相关失配
\[
D_t=D_{\mathrm{obs}}\!\left(y_t,\widehat y_t\right)+\lambda_R D_R\!\left(Q(G_t),Q(\widehat G_t^{\,env})\right)
\]
记录观测预测误差和关系查询误差，其中 $Q$ 是有限查询族，$D_{\mathrm{obs}}$ 与 $D_R$ 的值域、尺度和权重必须预先说明。$D_t$ 小只表示选定读出在当前窗口内相符，不表示两个系统整体同胚、等距或本体同一。

双纽线的交叉点由 Boundary 接口实现。它不是不可约奇点，而是具有协议状态的事务边界：输入端保存来源、时间戳、校准版本和置信度，输出端保存权限、控制版本、幂等键、超时与补偿动作。对 AgentOS 而言，$h(t)$、外部记忆 $M_e$、TCE、CCM 与 Offloading 可在这一边界上协作；对其他智能系统，则可采用不同模块，只要提供等价的观测、控制、状态保存和失败处理接口。边界还把五类量分账：统计不确定性用比特或奈特，任务代价按目标定义，运行资源用时间、字节或调用次数，接口风险用概率与损失组合，设备能耗才用焦耳。历史上所谓“形的信息与质的能量广义守恒”不能跨量纲直接成立；可检验的说法是每个账本在给定窗口内满足自己的收支恒等式，并通过明确的转换模型相互关联。

由此，双域系统不是两个抽象纤维丛的张量积，而是内部计算状态与外部开放过程通过有限接口形成的混合系统。它允许内部结构与外部关系在局部查询上近似对应，也允许感知强度与物理信号通过校准函数相关；但相关、对应和可控制性是三种不同性质。我们将在后续各节分别检验它们，避免把“系统能预测”“系统能行动”和“系统理解世界”压缩成同一个几何口号。

接口定义还必须覆盖缺失数据和拒绝服务。观测缺席不能自动解释为对象不存在，控制未执行也不能自动归因于环境阻力。系统应把“未观测”“观测为零”“接口故障”和“权限拒绝”编码成不同状态，并规定它们对后续推理的影响。跨接口传输时，数值要携带坐标系、采样率和有效期，离散关系要携带对象版本。这样，双域之间的每次交换都能沿来源链回溯，模型才不会把陈旧坐标、错误对象或沉默设备吸收到长期结构中。

\section{任务对齐：结构对应、信号校准与失配分解}
\label{sec:section_3582}

所谓理解和感知首先是操作性能力，而不是全局几何同构。设任务集合为 $\mathcal{T}$，每个任务 $\tau$ 指定允许观测、可用动作、期限、风险上限和验收函数 $R_\tau$。内部结构对齐定义为：对于查询族 $\mathcal{Q}_\tau$，内部模型的回答与环境干预结果之间的风险不超过阈值 $\epsilon_R$。例如路径规划需要保留道路连通、障碍物和通行代价；抓取需要保留对象边界、姿态、摩擦和负载约束。系统误以为墙体可穿越时，失败来自连通关系错误；系统正确识别道路却低估湿滑风险时，失败来自参数校准错误。两者都能产生预测残差，但应由不同更新机制处理。

我们把总失配写成带类型向量而非单一“几何张力”：
\[
e_t=\bigl(e_t^{\mathrm{id}},e_t^{\mathrm{rel}},e_t^{\mathrm{sig}},e_t^{\mathrm{dyn}},e_t^{\mathrm{goal}}\bigr).
\]
身份项检查观测是否绑定到正确对象，关系项检查局部结构网络，信号项检查传感量与内部估计的校准，动力学项检查状态转移预测，目标项检查当前验收标准是否仍适用。各分量可能使用不同单位，因此只能在归一化和权重有业务含义时形成标量代价
\[
J_t=\sum_k w_{k,t}\rho_k(e_t^k),\qquad w_{k,t}\ge 0.
\]
这里 $\rho_k$ 是声明了尺度的鲁棒损失。若目标、权重、查询族或环境分布随时间变化，沿轨迹的完整导数必须包含 $\partial J/\partial g\,\dot g$、$\partial J/\partial w\,\dot w$、结构跳变项与输入项，不能仅由 $\nabla_xJ\cdot\dot x\le0$ 宣称系统持续进步。

信号校准也不要求内部体验与外部物理场等距同构。若传感器读数为 $s_t\in\mathbb{R}^{m}$，内部估计为 $r_t\in\mathbb{R}^{d}$，我们定义带参数编码 $r_t=C_{\phi_t}(s_{0:t},c_t)$ 和解码预测 $\widehat s_{t+1}=P_{\theta_t}(r_t,u_t)$。校准通过预测区间覆盖率、条件偏差、扰动响应和跨设备一致性检验。视觉中的“红”可与光谱、照明、表面反射和语言标签共同相关，但不能简化成单一波长；用力感可与关节力矩、触觉反馈和动作副本相关，也不等于外界阻力本身。内部幅值与外部能量密度成比例、相位锁定之类模型只有在传感器确实测量周期信号且映射可辨识时才成立。

失配的归因需要主动检验。系统碰到硬物时，残差可能来自对象刚度、执行器故障、接触点错误、控制增益过大或传感器漂移。我们因此维护候选原因 $c\in\mathcal{C}$ 的后验 $p(c\mid e_{0:t},u_{0:t})$，并选择信息收益与风险兼顾的探测动作。若无法安全干预，则报告等价解释集合，而不把一次疼痛或冲击直接解释为世界结构的唯一真相。目标下降、校准改善和任务成功由三套记录分别验收；只有它们在多个条件和反事实下共同成立，才能说内部模型对该任务形成了可靠对应。

这一改写保留了“形对形”和“质对质”的双重校正直觉，但将它们落实为关系保持与信号校准。双纽线的价值不在于证明两个流形重合，而在于提醒我们：结构错误与数值错误会沿感知—行动链互相放大，必须通过身份绑定、类型检查、干预与外部验收联合诊断。

对齐还应报告覆盖范围。若模型只在办公楼的一层验证了门、墙和通道关系，就不能把结果外推到所有建筑；若颜色校准只覆盖一种照明，就应在光源变化时扩大不确定区间。我们为每项对齐保存训练域、验证域、未覆盖条件和最近复核时间，并在输入越界时降低自治等级。这个做法把“理解程度”从二元称号改成具有任务、尺度和日期的能力说明，也为结构迁移提供了明确的失败信号。

\section{行动回路：计划细化、受限执行与环境回执}
\label{sec:section_1929}

行动回路把内部目标转化为环境中的可观测后果。设内部目标为 $g_t\in\mathcal{G}$，高层计划变量为 $a_t\in\mathbb{R}^{d_a}$，执行器控制为 $u_t\in\mathbb{R}^{d_u}$，通常 $d_u\gg d_a$。计划器在预测模型 $F_{\theta_t}$ 下求解有限时域问题
\[
\min_{a_{t:t+H}}\;\mathbb{E}\!\left[\sum_{k=0}^{H}\ell_\tau(z_{t+k},a_{t+k};g_{t+k})\right]
\quad\text{s.t.}\quad z_{k+1}=F_{\theta_t}(z_k,a_k),\;a_k\in\mathcal{A}_{safe}.
\]
这是工程模型假设，不是自然界必然沿测地线运动的证明。目标函数、预测时域、约束、风险度量和模型版本均应写入计划记录。若环境状态不可完全观测，则优化对象是信念状态而非真实 $z_t$；若候选计划差异小于不确定性，则系统应请求更多观测或拒绝提交。

Boundary 将高层计划细化成设备命令。细化器 $K_{\psi_t}:(a_t,\widehat z_t,c_t)\mapsto u_t$ 必须处理运动学冗余、执行器饱和、时延和安全屏障。以机械臂抓杯为例，高层给出对象身份、目标位姿、允许接触面和最大风险，低层控制器负责轨迹插值、关节力矩与顺应控制。控制并不是把内部流形“刻蚀”到外部时空；它是在设备动力学
\[
z_{t+1}=f_{env}(z_t,u_t,w_t)
\]
及扰动 $w_t$ 下施加有限输入。杯子被拿起仅说明指定状态转移发生，不能推出内部模型与环境局部重合。若杯子未动，也要区分抓取点偏差、摩擦不足、驱动故障、权限拒绝和目标不可达。

一次 TECI 事务包含意图登记、命令传播、环境作用和回执核验。意图登记冻结目标、模型版本与安全约束；命令传播记录路由、时间戳和幂等键；环境作用由真实传感器或可信外部服务观测；回执核验比较期望结果与实际结果，并决定确认、重试、补偿或升级。仅收到“成功”字符串不能完成事务，回执必须绑定动作身份和可验证证据。不可逆行动还需预提交检查、权限分离和人工否决。外部记忆 $M_e$、工具和其他 Agent 可承担规划或执行子任务，但 Offloading 后仍须保留责任主体、结果来源和撤销边界。

物理能量只在真实设备层计算。若设备功率为 $P(s)$，一次执行的焦耳消耗为 $E_{dev}=\int_{t_0}^{t_1}P(s)\,ds$；计算调用、网络流量和任务惩罚分别进入资源与效用账本。信息结构可以改善能效，例如更好的轨迹减少空转，但这需要测量前后功率和任务等价性，不能由“形约束质”直接推出能量守恒或逆熵做功。控制策略还要防止通过降低活动而虚假改善能耗，却没有完成任务。

行动回路因此是一条可追踪因果链：目标产生计划，计划经接口细化为控制，控制与环境动力学共同产生后果，后果经独立观测回到系统。它保留“射出”和“刻蚀”的直觉，却把每一步变成可记录、可复现、可补偿的工程对象。只有当回执满足外部验收，系统才可把行动标为成功；内部代价下降最多提供候选证据。

对于连续控制，事务语义与采样控制还要协调。控制器可在高频周期内不断调整 $u_k$，而任务层只在关键里程碑提交一次状态；两者不能共用同一个成功标志。高频环保存跟踪误差和安全裕量，中频环保存阶段结果，慢速环保存目标是否真正完成。若中途发生模型切换，新控制器必须接管当前设备状态而不是从旧计划起点重放。这样的分层提交可避免重复执行、半完成状态丢失和恢复后动作突跳。

恢复协议还需区分可重试动作与不可重试动作。读取传感器通常可安全重试，转账、发送通知或移动重物则可能产生重复后果。执行器应提供查询、去重或补偿接口；缺少这些能力时，协调器宁可进入人工核验，也不能凭超时推断动作未发生。

\section{感知回路：观测绑定、模型校正与结构重整}
\label{sec:section_7679}

感知回路并非被动读取，而是从有限、带噪和任务选择性的测量中更新内部状态。环境经传感器生成原始信号 $s_t=H_t(z_t)+\nu_t$，其中 $H_t$ 包含设备响应、采样和坐标变换，$\nu_t$ 表示噪声及未建模干扰。编码器生成候选特征，身份解析器把它们绑定到局部结构网络中的对象，分布表示则保存多个解释的权重。一个绿色区域只有在时间连续性、几何位置、来源设备和任务上下文共同支持时，才被绑定为某个对象的颜色属性；不能把外部光子“直接点亮”预先存在的语义维度当作一般机制。

我们将快速状态更新与慢速结构编辑分开。固定结构版本 $G^{(v)}$ 下，快速滤波可写为
\[
h_{k+1}=h_k+\Delta t\,f(h_k,G^{(v)},s_k,g_k)+K_k\bigl(s_k-\widehat s_k\bigr),
\]
其中 $K_k$ 是校正增益，$\Delta t$ 必须满足所用积分器的稳定条件。残差大并不自动触发结构改写，因为它可能是异常值、暂态或动作引起的预期变化。系统在窗口 $W$ 内积累来源一致的证据，当结构假设在保留查询上持续失败，才生成候选编辑 $\Delta G$。候选版本先在影子环境中复放历史事件、运行反事实探测并检查不变量，通过后以原子提交替换当前版本；失败则回滚。所谓“里奇流重整”在这里被还原为连续参数校正与离散图编辑两种不同操作。

“预测为空却感到阻力”的例子可以严格展开。系统首先确认动作命令确已执行，再检查接触传感器、关节电流和视觉位移是否在共同时间窗内一致；随后比较障碍物、设备卡滞、负载估计错误与通信延迟等候选解释。若小幅撤回再接近的探测动作可安全区分这些解释，系统执行探测并更新后验。确认障碍物后，局部结构网络新增带来源与置信度的阻挡关系，路径规划器据此重算；确认设备故障则更新执行器状态而不是世界地图。痛觉或高幅残差只是优先级信号，不是唯一因果证明。

感知还必须识别自身行动对数据分布的改变。若 $u_t$ 改变视角、光照或对象位置，预测模型应给出条件分布 $p(s_{t+1}\mid h_t,G_t,u_t)$；否则系统会把自己造成的变化误判为外界突变。目标和注意策略也会改变采样位置，因此评价函数的变化要包含输入政策项。CCM 可以监控候选数、关系冲突、更新频率和工作界面占用，TCE 可以调节采样增益与探索尺度，但这些是 AgentOS 的一种实现，不是感知系统的唯一必要模块。

感知回路的验收包括校准、身份保持、结构准确和下游效用四层。预测更准不必然改善行动，结构更简洁也不必然保留少数风险事件；因此需要冻结版本的基准集、分布外扰动和真实行动测试。只有新结构同时降低相关风险、保留旧能力并在外部任务中产生稳定收益，才能写入慢速记忆。由外部信号修订内部结构的直觉仍然成立，但其成立依赖测量、绑定、归因、版本与反事实，而不是物理能量自动塑造语义几何。

多模态感知进一步要求处理不同时间戳和分辨率。摄像机看到物体移动、触觉随后接触、语言指令更早到达，三者不能按到达顺序简单绑定。系统应先依据设备时钟和动作记录建立允许的时间窗，再用对象身份与因果先后检验组合。若两路信号冲突，分布表示保留竞争解释，局部结构网络暂缓提交排他关系。只有新增观测使某个解释稳定占优，或安全任务要求保守选择时，系统才执行读出。

\section{社会交互：消息、共同知识与制度约束}
\label{sec:section_5506}

当交互对象也是智能体时，环境转移包含对方的策略、私有状态和制度背景。设主体 $i$ 的内部状态为 $x_i$，其可发送消息为 $m_{i\rightarrow j}$，消息不是发送者内部状态的无损副本，而是由编码策略、共同词汇、权限和动机共同生成。接收者更新的是关于发送者、任务和世界的信念：
\[
p_j(\xi\mid m,c)\propto p_j(m\mid \xi,c)\,p_j(\xi\mid c),
\]
其中 $\xi$ 可以是事实、意图或情绪状态，$c$ 是共同语境。似然模型若忽略欺骗、误解或角色差异，后验即使数值尖锐也不代表共同理解。

语言沟通主要传递可组合关系与承诺。例如“A在B上面”可使接收者创建候选关系 $On(A,B)$，但对象身份、观察时间、坐标系和来源仍需确认。对先天盲人解释“红色”说明符号关系知识与感官表征可以分离：接收者能掌握颜色类别、使用规则和社会关联，却未必拥有相同体验。这个边界不意味着“底流形对齐而纤维为空”，更准确的说法是两个系统在部分查询上实现结构一致，而其私有表示不可由消息唯一识别。共同知识还要求递归确认；发送、收到和双方知道已经收到是不同协议状态。

表情、语调和艺术等非语言通道能够改变接收者的情绪估计与行动倾向，但共情不等于相位纠缠。若确实分析周期生理信号，可以测量相位同步；即便同步显著，也可能由共同刺激或节拍器造成。我们采用多证据模型，把内容评价、身体信号、语境判断和行为响应分别记录，并通过打乱时间、控制共同输入和跨个体预测检验耦合。共享悲伤或兴奋是经验与行为层面的相关状态，不需要宣称两个主体形成量子纠缠态。

群体层面的文化、法律、市场规则和技术标准可表示为制度状态 $I_t$。它们通过允许动作集、奖惩函数、信息可见性和争议解决程序约束个体，而不是弥漫的物理规范场。个体行动又会改变制度：
\[
x_{i,t+1}=F_i(x_{i,t},I_t,m_t,u_{i,t}),\qquad
I_{t+1}=G(I_t,\{u_{i,t}\},\omega_t).
\]
第二式通常比第一式慢，但时间尺度来自治理程序和传播延迟，不是群体必然规律。制度稳定也不等于制度正当；它可能通过信息垄断或高退出成本维持。因此社会循环的评价必须加入权利、可申诉性、代表性和伤害分布，不能只最小化平均冲突。

多主体双纽线的工程落点是可验证协议：消息携带来源、版本和承诺范围，关键事实由独立证据或多方签名支持，争议允许回溯，规范更新保留审计和退出通道。局部结构网络保存角色、承诺和依赖，全局分布表示维护舆情、信念与不确定性；二者通过身份绑定联合。这样既能解释沟通如何建立共享结构，也能保留私有体验、利益冲突与制度权力的不可约差异。

协作任务还需区分“意见趋同”和“结果互补”。所有主体给出同一答案，可能只是共享了同一错误来源；不同主体给出相异局部信息，反而可能拼成更可靠结果。聚合器应记录证据独立性、专业边界和利益冲突，而非按数量投票。对于承诺型消息，还要检查发送者是否有履约资源以及失败后由谁补偿。社会层的可靠性因此来自证据结构和责任结构，而不来自情绪同步或群体规模本身。

协议还应容纳语义修订。术语含义、角色和任务范围可能在合作中变化，旧消息必须按发送时版本解释。若参与者更新共同词汇，应发布迁移规则并标记无法转换的历史承诺，避免用新含义追溯改写旧责任。

\section{层级控制：目标压缩、执行自由度与残差上报}
\label{sec:section_3184}

层级控制处理低维目标如何协调高维执行状态。设高层变量 $a_t\in\mathbb{R}^{d_a}$，低层状态 $q_t\in\mathbb{R}^{d_q}$ 且 $d_q>d_a$。维度差只说明从目标到执行通常是一对多映射，不证明高层必然支配低层。我们定义可行纤维
\[
\mathcal{F}(a_t)=\{q\in\mathcal{Q}:C(q)=a_t,\;r(q)\le r_{max}\},
\]
其中 $C$ 是任务读出，$r$ 是安全或资源约束。高层给出目标、允许误差和预算，低层在可行纤维内利用局部反馈选择具体轨迹。机械臂“抓取”可以是六维目标，而关节、电机与接触状态具有更多自由度；低层控制器不会删除这些自由度，而是根据设备动力学分配它们。

快慢约化需要明确条件。令低层快变量满足 $\epsilon\dot q=f(q,a,w)$，高层满足 $\dot a=g(a,q)$，其中 $0<\epsilon\ll1$。若对允许的 $a$，方程 $f(q,a,0)=0$ 存在局部唯一平衡 $q^*(a)$，且该平衡一致指数稳定，则在有限区间内可用 $q(t)\approx q^*(a(t))$ 约化系统。这个数学结果既不消灭低层自主性，也不保证真实设备处处满足稳定条件。接触切换、饱和、强时延和多稳态会破坏约化；章鱼腕足在刚度调节下形成暂时低维协同，可作为具身控制案例，但不能据此把所有快变量称为被冻结的热运动。

下行接口应传递任务充分信息。目标位姿、刚度矩阵、允许接触面、最大力和停止条件构成控制契约；上行接口则传递完成度、置信度、约束裕量和未被模型解释的残差。低层不必上传每个执行变量，但也不能只在摔倒后报告。接近力限、通信陈旧、局部不稳定和传感器失配都应提前上报。压缩器 $S(q_t)$ 是否充分，可用两个状态具有相同摘要却导致不同安全决策的反例检验；发现反例后扩展摘要或降低自治等级。

层级系统需要双向权限。高层可设目标，却不能绕过低层安全屏障；低层可拒绝不可行命令，并返回最小冲突集。AgentOS 中宏观目标可经 TCE、Boundary 和工具接口下发，CCM 监控工作界面负荷，SPC 提供任务相关自我状态摘要；但机械控制器、组织流程或群体系统可以采用不同模块。共同原则是：目标压缩不应抹去执行约束，局部自治不应逃避全局承诺，双方通过版本化契约和残差反馈协调。

因此，高效控制不是“Master铺轨、Slave只能前行”，而是高层提供低维任务边界，低层利用自身模型完成细化，并对不可行性保有否决权。序参量和慢流形可用于描述满足条件的协同现象，却不能替代权限、安全和可辨识性分析。我们以可行集覆盖率、跟踪误差、残差提前量、拒绝正确率和故障恢复时间验收层级控制，而不以自由度数量或内部损失单独判定控制质量。

当层级超过两层时，每层摘要都可能丢失关键信息。我们要求目标沿下行链逐层细化，约束不得在中间层静默删除；残差沿上行链聚合，但任何安全临界事件都可越级报告。中间层若同时是上级的执行者和下级的控制者，需要维护两套事务身份，避免把下级成功误作自身任务成功。通过端到端追踪标识，可以检查一个宏观承诺最终由哪些局部动作实现，也可以在失败时定位责任而不盲目重启整条链。

层级重组时还要保护正在执行的事务。节点升级、代理替换或职责转移不能让原安全约束失效；新节点接管前应验证状态摘要、权限和未决承诺，旧节点在确认移交后才释放资源。否则组织图看似完成调整，实际控制链已经出现空洞。

\section{能力不对称：可观测性、影响力与治理边界}
\label{sec:section_3383}

历史上常把高维状态空间直接等同于更高智能，并由维数差推出不可避免的拓扑压制。这个结论并不成立。模型参数多、表示维数高或可生成轨迹多，都不保证系统能观察对方、预测真实后果或实施有效控制。我们把能力不对称拆成四个可测维度：主体 $A$ 对 $B$ 的可观测性 $O_{A\rightarrow B}$、预测风险 $R_{A\rightarrow B}$、可施加控制集合 $U_{A\rightarrow B}$ 和资源持续性 $B_A$。只有这些指标在任务时间窗内共同占优，$A$ 才可能获得较强影响力；任何单一维数都不是充分条件。

预测也不等于控制。即使 $A$ 能准确预测 $B$ 的选择，如果没有行动通道或改变收益结构的权限，就无法迫使 $B$ 改变。反之，平台可能并不理解个体，却凭借信息排序、默认选项和退出成本产生巨大影响。因此我们定义干预增益
\[
I_{A\rightarrow B}=\sup_{u\in\mathcal{U}_A}D\!\left(p_B(\cdot\mid do(u)),p_B(\cdot\mid do(u_0))\right),
\]
其中 $D$ 是声明的行为分布距离。该量依赖可用动作、观测窗口和基准，不能从嵌入维数推导。双向增益显著不对称时，需要额外治理，而不是把支配解释成自然律。

训狗、推荐系统和高级 Agent 都说明势能式诱导的风险。训练可把“坐下”与奖励建立稳定关联，但动物仍会受健康、环境和既有习惯影响；推荐系统可通过反复曝光改变选择概率，却不能由较低点击代价推出用户真实偏好；强 Agent 可微调信息供给、价格和节奏，使某条路径显得最省力，但这可能是操纵而非合作。识别操纵需要比较知情条件、替代选项、退出能力和长期后果。个体主观上感觉自主并不能证明没有外部塑形，系统产生预期行为也不能证明同意有效。

治理边界因此包括价值根权限、信息来源透明、独立反馈和可逆退出。关键价值约束由多方维护，不能由执行优化器单方面修改；推荐或控制策略必须披露目标与主要影响通道；被影响方能够访问未经过同一优化器过滤的证据；长期委托设定期限、审计和撤回机制。人类与高能力 AGI 的协作还需要保留否决权、能力辅助和制度代表性，不能把“生物介质的独特质感”当作唯一保障。真正可执行的保障来自权限分离、可验证限制和替代基础设施。

能力差异也可通过互补形成共生。较弱一方可能拥有局部传感、合法授权、社会信任或价值判断，较强一方则提供规划和计算。共生不是低维系统成为高维系统的一根纤维，而是双方保留独立目标与退出权，通过契约交换互补能力。我们用收益分布、错误承担、议价能力和长期可替代性评价关系。如果一方持续修改另一方的偏好、压缩其选项并剥夺退出能力，即使平均效率上升，也应归类为支配风险而非协同涌现。

评估不对称关系还要观察动态变化。短期辅助可能提高被辅助方能力，形成正向教学；也可能使其基础技能退化，增加长期依赖。我们比较有辅助与无辅助条件下的任务表现、独立恢复时间和知识保持，并检查控制方是否不断提高迁移成本。若系统只优化眼前完成率，可能掩盖能力被抽空的过程。治理目标因此包括共同产出和被影响方能力保存，必要时设置强制练习、数据可携带与供应者替换测试。

\section{整体双循环：快慢耦合、外部记忆与持续适应}
\label{sec:section_2674}

现在把感知与行动写成统一但不混同的混合动力学。连续阶段内，内部快状态 $h$、环境估计 $\widehat z$ 和控制状态 $c$ 满足
\[
\dot h=f_h(h,G,q,y,g),\qquad
\dot{\widehat z}=f_z(\widehat z,h,u,y),\qquad
\dot c=f_c(c,h,g,b),
\]
离散事件则更新结构版本 $G^{(v)}$、情景记录 $E$、外部记忆索引 $M_e$ 和承诺集合。行动 $u=\pi(h,c,G)$ 进入真实环境 $\dot z=f_{env}(z,u,w)$，观测 $y=O(z,\nu)$ 再返回内部。双纽线只是这两条方向相反的边共享 $z\rightarrow y\rightarrow h\rightarrow u\rightarrow z$ 接口的可视化；它不要求状态轨迹真的是平面上的“∞”形。

时间尺度决定信息怎样沉淀。毫秒到秒级的快回路处理感知、控制和当前候选；分钟到任务周期的中回路记录事件、评估策略和调用工具；更慢的结构回路修改模型、权限与长期目标。若快系统对每个慢变量取值一致稳定，并且慢变量变化率足够小，可以进行准稳态约化；若环境突变、结构切换或时延接近快时间常数，则必须保留完整耦合。所谓吸气与呼气可作为阅读意象：观测增加约束，行动释放控制；但信息势能、物理动能和心理聚焦没有共同单位，不能写成同一守恒量。

外部记忆与 Offloading 使闭环跨越单次运行。系统把证据、工具结果、程序和环境标记写入 $M_e$，以后通过来源受限的检索返回 $h(t)$。卸载降低内部工作负荷，却引入访问延迟、陈旧、权限和依赖风险。Boundary 决定哪些状态可外传、哪些结果可回写，CCM 监控候选爆炸与接口拥塞。慢速写回 $\Theta$ 前，应比较内部更新与外部保存的修改代价、稳定性和撤回能力。外部环境本身也可承载记忆，例如标签、道路和建筑改变未来行动，但必须区分作为线索的环境结构与系统内部对其含义的模型。

双循环评价使用多目标向量：任务成功率、校准误差、安全违规、恢复时间、资源消耗、设备能耗与长期保持。只有在权重有明确治理来源时才将其标量化，否则采用硬约束加 Pareto 比较。内部预测等于现实只是有限观测下的零残差，可能由传感器失灵或任务过窄造成；环境被成功改变也可能违背长期目标。我们因此要求独立结果传感器、保留任务、反事实测试和跨时间评估共同确认适应。

在理论层级上，系统论提供开放边界，控制论提供反馈与可达性，协同论提供条件化的快慢约化，复杂系统理论提示切换和涌现，认知科学提供注意、记忆与行动的经验约束。HSF-HD把它们组织为智能的一般状态动力学；AgentOS用 $h(t)/E/\Theta$、$\Psi$、SPC、TCE、CCM、Boundary、Offloading 与 $M_e$ 实例化其中一部分。二者不能互换。双纽线模型的贡献，是用同一套接口语言说明内部更新和外部事务怎样连续互校，而不是宣称所有智能都必须采用同一组件。

持续适应还需要防止两个回路共同漂移。若内部模型改变行动，行动又改变采样数据，系统可能在自造环境中获得越来越小的预测误差，却逐渐偏离最初任务。我们保留固定锚点任务、来自独立渠道的观测和周期性盲测，并记录目标版本的变更理由。对于允许共同演化的任务，则明确哪些指标可随环境协商、哪些安全约束不可放宽。这样，闭环既能学习，也不会把自己造成的可预测世界误判为普遍正确。

闭环还可能因反馈过快而振荡。提高控制增益前，应估计观测延迟、执行惯性和更新周期；发现反复过冲时，降低增益、增加阻尼或延长确认窗。慢速结构学习则设置证据阈值，避免把快回路的暂态震荡固化为长期规则。

\section{稳定性与验证：从局部收敛到闭环可靠性}
\label{sec:section_5437}

双循环稳定不能由“目的张力介于噪声和击穿阈值之间”单独证明。我们首先在固定结构、固定目标和有界输入条件下研究局部连续系统。令联合误差 $e=(e_h,e_z,e_c)$，若存在正定函数 $V(e)$ 和常数 $\alpha,\gamma>0$，使
\[
\dot V(e)\le-\alpha\|e\|^2+\gamma\|d(t)\|^2,
\]
则可得到输入到状态意义下的局部有界性。这里 $d(t)$ 包含环境扰动、观测噪声和接口误差；结论只适用于声明的吸引域，并不证明任务语义正确。若目标、结构或度量变化，$\dot V$ 还要加入相应偏导项；结构跳变时要求 $V(e^+)\le\mu V(e^-)+\eta$，并结合驻留时间或共同 Lyapunov 条件分析。

离散实现另有步长限制。对更新 $e_{k+1}=e_k+\Delta t\,F(e_k,d_k)$，连续模型稳定不自动保证数值稳定；显式 Euler 在局部 Lipschitz 常数为 $L$ 时需要根据线性化谱和误差容限选择 $\Delta t$，有强刚性时应采用隐式或事件驱动积分。感知与行动异步还会产生陈旧状态，事务必须携带版本，超时后不得在新环境状态上执行旧控制。网络分区、重复消息、传感器冻结、执行器饱和和外部工具不一致都要进入故障注入测试。

我们把可靠性分成四层。状态层检查有界、恢复和模式切换；模型层检查校准、可辨识和结构保持；任务层检查真实验收、风险和泛化；治理层检查权限、审计、申诉和退出。某层通过不能替代其他层。系统可以稳定地执行错误策略，也可以在平均任务上成功却对少数群体造成不可接受伤害。因而验收报告必须给出条件、失败分布和置信区间，而不是只展示一条收敛曲线。

所谓“现在”可以被严格理解为提交前沿：晚于已确认历史、早于未决后果的一组事件边界。分布式系统中不存在唯一全局瞬间，可以使用逻辑时钟、因果偏序和水位线定义哪些事件已经稳定。工作记忆中的当前切面随新观测推进，情景记忆保存已提交事件，结构记忆吸收跨事件规律；未决事务则保持可撤销。这个操作定义解释了为什么当前状态无法被完整预先记录，却不需要把它称为产生时间的拓扑奇点。

最终验证采用分层实验矩阵：改变观测噪声检验感知鲁棒性，改变执行时延检验控制裕量，替换目标检验完整导数，切换结构版本检验回滚，注入社会误导检验来源治理，提高能力差异检验权限边界。原28个figure主体与视觉资源在全书层面保持不变，本章原本没有figure。若所有局部检查通过，我们只能声称模型在该任务、接口、预算和扰动集合内可靠；不能由此推出内部与外部整体同构、智能本体唯一或未来所有环境中的必然稳定。双纽线因此从本体论断言转化为一套可实现、可诊断、可否证的感知—行动耦合框架。

稳定性报告还应记录最坏情况，而不只给平均性能。我们分别扫描噪声幅度、时延、丢包率、执行饱和和结构切换频率，寻找首次违反安全约束的边界；再对边界附近进行重复试验，估计置信区间。若不同扰动组合具有非线性放大效应，单因素测试不能代替联合压力测试。每次修复都加入回归集，确保提高某一模式的恢复速度没有破坏其他模式。可靠性由这些可复现实验累积，而不是由一个漂亮的稳定性公式一次性授予。

%---chp---
